% Chapter 10: Accuracy, cost, limits and guarantees (practice author). PLAN.md section 1.4.
% Tables: T14 tab:accuracy-matrix (generated, table*), T15 tab:cost (generated), tab:ridders-statistics (values of
% DS9), tab:tests. Back-references to F14 fig:ridders-coverage and F15 fig:ridders-edge.
\chapter{Accuracy, cost, limits and guarantees}\label{ch:guarantees}

This chapter states what to expect of \code{nxx::deriv}, and on what evidence. It covers the accuracy per method,
order and type; the cost in evaluations; where the methods stop working; what the library guarantees, and under which
hypotheses; and how all of it is tested. Each statement is marked as measured, derived or a rule of thumb, and points
to its source.

\section{Accuracy}\label{sec:guarantees-accuracy}

\Cref{tab:accuracy-matrix} gives the measured relative error of seven methods at their default steps, for the first
to the fourth derivative, in three types. \code{BackwardRidders} is not among them, since it is the mirror of
\code{ForwardRidders}. The function is $\mathrm{e}^x$. All its derivatives equal the function, so the orders can be
compared. Each cell is the median over \num{\valAccuracyMatrixPoints} points $x = 0.5, 0.55, \ldots, 1$. These points
share the scale $s = 1$, and so they share the step. The median matters. At a single point, the rounding part of the
error is a matter of luck, by about a factor of ten. \code{Central<2>} for $f'$ shows it: at $x = 1$ its error is
\num{\valExpCentralTwoAtOneRelError}, against a median of \num{\valExpCentralTwoRelError} (measured).

\begin{table*}
\caption{Measured accuracy: the median relative error of the derivative of $\mathrm{e}^x$ over
\num{\valAccuracyMatrixPoints} points in $[0.5, 1]$, at the default steps, by method, derivative order N and type.
$^{\ast}$\,Every float derivative above the first warns; these values were computed with the warning silenced
(\cref{sec:api-warning}). 1.4e$-$5 stands for $1.4 \times 10^{-5}$.\label{tab:accuracy-matrix}}
\centering\footnotesize\setlength{\tabcolsep}{2pt}\sisetup{output-exponent-marker=\ensuremath{\mathrm{e}}}
\begin{tabular}{@{}lcccc@{\hspace{6pt}}cccc@{\hspace{6pt}}cccc@{}}
\toprule
 & \multicolumn{4}{c}{\code{float}} & \multicolumn{4}{c}{\code{double}} & \multicolumn{4}{c}{\code{cpp_bin_float_50}}\\
\cmidrule(lr){2-5}\cmidrule(lr){6-9}\cmidrule(l){10-13}
Method & $N = 1$ & 2 & 3 & 4 & $N = 1$ & 2 & 3 & 4 & $N = 1$ & 2 & 3 & 4\\
\midrule
\tablebody{accuracy-matrix}
\bottomrule
\end{tabular}
\end{table*}

The error model of \cref{ch:stencils} explains the table. A rule of accuracy order P at its default step has an error
of about $\epsmach^{P/(P+N)}$, relative (\cref{tab:default-steps}). Compare the double column with the model. The
central rules lie below the model by up to a factor of \num{\valModelRatioCentralBelowMax} for $f'$ to $f'''$, and
up to \num{\valModelRatioCentralFourthAbove} times above it for $f''''$. \code{Forward<2>} lies up to
\num{\valModelRatioForwardTwoAbove} times above it. \code{Forward<4>} lies \num{\valModelRatioForwardFourThird} and
\num{\valModelRatioForwardFourFourth} times above it for $f'''$ and $f''''$, because of its larger amplification
$\ampl$, which the model drops (measured). In round numbers, \code{Central<2>} in double gives
\num{\valHeadlineFirstDigits} correct digits for $f'$ and \num{\valHeadlineFourthDigits} for $f''''$. Ridders' method
gives \num{\valHeadlineRiddersFirstDigits} and \num{\valHeadlineRiddersFourthDigits}. \code{ForwardRidders} loses
more with each order, because its columns gain one order rather than two. On the random study of
\cref{sec:guarantees-ridders}, its median error is \num{\valForwardRiddersAccuracyFactorFirst},
\num{\valForwardRiddersAccuracyFactorSecond}, \num{\valForwardRiddersAccuracyFactorThird} and
\num{\valForwardRiddersAccuracyFactorFourth} times Ridders' for $f'$ to $f''''$ (measured). Higher orders pay most in
the wide type, where the step can grow far before truncation matters. The float columns show why every float
derivative above the first warns, whatever the method.

Two caveats apply to every figure. First, the model gives an absolute error, $\epsmach \abs{f} / h^N$ in its rounding
part. The relative figures presuppose that $\abs{f} \sim s^N \abs{f^{(N)}}$. That is true for $\mathrm{e}^x$ on
$[0.5, 1]$, and nearly so for the nitrogen $\cp$, where $\abs{\enth} / (s \abs{\cp}) = \num{\valCpRelativeFactor}$
(\cref{ch:rounding}). Near a zero of $f^{(N)}$, the relative error grows without bound, while the absolute error stays
as it is (\cref{sec:using-check}). Second, the default steps assume that $f$ varies on the scale $\max(\abs{x}, 1)$,
and that it is accurate to $\epsmach$. Where $f$ varies faster or is noisier, the figures do not hold until the step
is fitted (\cref{fig:step-scale,fig:noise}).

\section{Ridders' error estimate}\label{sec:guarantees-ridders}

Ridders' method returns an estimate of its error, not a bound. We measured the quality of the estimate on a random
study (\cref{fig:ridders-coverage}). The study takes \num{\valRiddersFunctions} smooth functions with closed-form
derivatives ($\mathrm{e}^x$, $\sin x$, $\log x$, $\sqrt{x}$, $1/x$, $1/(1 + 25x^2)$, $x^5$, $\mathrm{e}^{-x^2}$,
$\tanh x$), each at \num{\valRiddersPointsPerFunction} random points in $[0.2, 3]$. It works in double, with the
default initial steps and no noise. For $f'$, no singularity lies within the initial step. For $f''$ to $f''''$, the
initial step $s/4$ reaches past zero for $x$ below 0.25, and the span $2 h_0$ of $f'''$ and $f''''$ does so for $x$
below 0.5. These cases are the restarts of \cref{tab:ridders-statistics}. The table gives the results per derivative
order. The coverage is the fraction of cases whose estimate is at least the actual error. The ratio is the actual
error over the estimate, so a ratio below 1 is covered.

\begin{table}
\sidecaption{Ridders' method on the random study: \num{\valRiddersCases} cases per derivative order, measured. The
restarts counted are those after a value of $f$ that is not finite, counted from a log of the calls. A restart after
an untrusted result continues the halving of the step, and cannot be told apart.\label{tab:ridders-statistics}}
\centering\footnotesize
\begin{tabular}{@{}lcccccc@{}}
\toprule
Order & Coverage & Uncovered & \multicolumn{2}{c}{Ratio} & Evaluations & Restarts\\
\cmidrule(lr){4-5}
 & (\%) & & median & worst & min, median, max & \\
\midrule
$f'$ & \num{\valRiddersCoverageFirst} & \num{\valRiddersUncoveredFirst} & \num{\valRiddersMedianRatioFirst} &
  \num{\valRiddersWorstRatioFirst} & \num{\valRiddersEvalsMinFirst}, \num{\valRiddersEvalsMedianFirst},
  \num{\valRiddersEvalsMaxFirst} & \num{\valRiddersRestartedFirst}\\
$f''$ & \num{\valRiddersCoverageSecond} & \num{\valRiddersUncoveredSecond} & \num{\valRiddersMedianRatioSecond} &
  \num{\valRiddersWorstRatioSecond} & \num{\valRiddersEvalsMinSecond}, \num{\valRiddersEvalsMedianSecond},
  \num{\valRiddersEvalsMaxSecond} & \num{\valRiddersRestartedSecond}\\
$f'''$ & \num{\valRiddersCoverageThird} & \num{\valRiddersUncoveredThird} & \num{\valRiddersMedianRatioThird} &
  \num{\valRiddersWorstRatioThird} & \num{\valRiddersEvalsMinThird}, \num{\valRiddersEvalsMedianThird},
  \num{\valRiddersEvalsMaxThird} & \num{\valRiddersRestartedThird}\\
$f''''$ & \num{\valRiddersCoverageFourth} & \num{\valRiddersUncoveredFourth} & \num{\valRiddersMedianRatioFourth} &
  \num{\valRiddersWorstRatioFourth} & \num{\valRiddersEvalsMinFourth}, \num{\valRiddersEvalsMedianFourth},
  \num{\valRiddersEvalsMaxFourth} & \num{\valRiddersRestartedFourth}\\
\bottomrule
\end{tabular}
\end{table}

The estimate is built so that it errs on the safe side. An extrapolated entry's estimate is its change from its parent
in the tableau. But it is never less than the rounding floor of its level; for the one-sided forms, it is never less
than the floor carried through the tableau. The result's \code{abs_error} is also at least its distance to the latest
diagonal entry. The method stops when an entry equals its parent, after two levels that improve nothing, or at the
resolution of $x$. It restarts at half the smallest step that it tried, in two cases: after a step at which $f$ is not
finite, and after an untrusted result. An untrusted result is one whose estimate is at least a quarter of the value
and more than 64 times the rounding floor. In all, the method attempts at most \num{\valRiddersLevelCap} levels
(\cref{tab:ridders-constants}, \cref{sb:floor-coverage}). The library's tests accept a result whose actual error is at
most six times its estimate, and whose estimate is within twice that of an independent prototype for the same case.

We measured the one-sided \code{ForwardRidders} on the same cases, with its default initial step $s/32$. Its estimate
covered \SI{\valForwardRiddersCoverageFirst}{\percent}, \SI{\valForwardRiddersCoverageSecond}{\percent},
\SI{\valForwardRiddersCoverageThird}{\percent} and \SI{\valForwardRiddersCoverageFourth}{\percent} of them for $f'$
to $f''''$. The worst ratios were \num{\valForwardRiddersWorstRatioFirst}, \num{\valForwardRiddersWorstRatioSecond},
\num{\valForwardRiddersWorstRatioThird} and \num{\valForwardRiddersWorstRatioFourth}. The median errors were
\num{\valForwardRiddersMedianRelErrorFirst}, \num{\valForwardRiddersMedianRelErrorSecond},
\num{\valForwardRiddersMedianRelErrorThird} and \num{\valForwardRiddersMedianRelErrorFourth} (measured;
\cref{sec:ridders-onesided}). A pole ahead of $x$, inside the span that the forward rule samples, lies outside that
study. A separate study covers it: $1/(t - a)$ at 0, with \num{\valPoleAheadCases} poles from
$2^{-\num{\valPoleAheadClosestExponent}}$ to $2^{-\num{\valPoleAheadFarthestExponent}}$ ahead of $x$. For $f'$ and
$f''$, the estimate always covered the error. For $f'''$, it missed \num{\valPoleAheadUncoveredThird} times, with an
error at worst \num{\valPoleAheadWorstRatioThird} times the estimate. For $f''''$, it missed
\num{\valPoleAheadUncoveredFourth} times. Of these, \num{\valPoleAheadSilentFourth} were neither covered nor flagged
by a large estimate. They were \num{\valPoleAheadSilentRelErrorMinFourth} to
\num{\valPoleAheadSilentRelErrorMaxFourth} times off, with an error up to \num{\valPoleAheadWorstRatioFourth} times
the estimate: within the tests' factor of six. The same poles behind $x$ gave \num{\valPoleBehindUncovered} uncovered
results (measured). Near a known singularity ahead of $x$, repeat a one-sided $f''''$ at a second initial step.

The estimate assumes that $f$ is accurate to about $\epsmach$. We added a relative noise of $10^{-10}$, and then of
$10^{-8}$, to $\mathrm{e}^x$. Over \num{\valNoiseRiddersCases} points, the coverage then fell to
\SI{\valNoiseTenRiddersCoverage}{\percent} and \SI{\valNoiseEightRiddersCoverage}{\percent}, with worst ratios of
\num{\valNoiseTenRiddersWorstRatio} and \num{\valNoiseEightRiddersWorstRatio} (measured, \cref{sec:rounding-noise}).
Nor does the estimate see a jump inside its span, except as a large estimate (\cref{fig:kink}). And it does not see a
kink at $x$ at all (\cref{sec:ridders-trust}).

\section{Cost}\label{sec:guarantees-cost}

The cost of a derivative is counted in evaluations of $f$. These dominate whenever $f$ is more than a few arithmetic
operations. The library's own work is one weighted sum per stencil, and one tableau row per level of Ridders' method.
\Cref{tab:cost} gives the counts. For a stencil they are exact: one evaluation per node. For Ridders' method and
\code{ForwardRidders} they depend on $f$ and $x$. For these the table gives the median and the range of the random
study, and the cap, which no call exceeds. \code{BackwardRidders} costs what \code{ForwardRidders} costs on the
mirrored function. For $f'$ in double, Ridders' method spends \num{\valRiddersEvalsMedianFirst} evaluations (the
median), against a stencil's 2 to 6. That buys about three more digits than \code{Central<2>}
(\num{\valHeadlineRiddersFirstDigits} against \num{\valHeadlineFirstDigits}): as many as \code{Central<4>} gives, and
no more than \code{Central<6>} (\cref{tab:accuracy-matrix}). It also buys an estimate of the error.

\begin{table}
\sidecaption{Evaluations of $f$ per derivative, by method and derivative order N. For the stencils the counts are
exact. For \code{Ridders} and \code{ForwardRidders} they are measured on the random study (the median, with the
range), and the cap is given. \code{BackwardRidders} costs as \code{ForwardRidders} does.\label{tab:cost}}
\centering\footnotesize
\begin{tabular}{@{}>{\raggedright}p{36mm}cccc@{}}
\toprule
Method & $N = 1$ & 2 & 3 & 4\\
\midrule
\tablebody{cost}
\bottomrule
\end{tabular}
\end{table}

\section{Limits}\label{sec:guarantees-limits}

\paragraph{Orders.} The derivative order is 1 to 4. The caps on P are design decisions
(\code{docs/deriv/PHASE2_API.md}). \code{Central<P>} stops at $P = 10$. By the model's exponent, a \code{Central<12>}
would cut the error of $f'$ by a factor of about \num{\valCapGainFactorCentralTwelveFirstDouble} in double and
\num{\valCapGainFactorCentralTwelveFirstMultiprecision} in a 50-digit type
(\num{\valCapGainFactorCentralTwelveFourthDouble} and \num{\valCapGainFactorCentralTwelveFourthMultiprecision} for
$f''''$), on a still wider stencil. The one-sided rules stop at $P = 6$. There the rounding amplification $\ampl$ is
already \num{\valAmplForwardSix} for $f'$ and \num{\valAmplForwardSixFourth} for $f''''$
(\cref{tab:weights-forward-first,tab:weights-forward-fourth}), and it grows about as fast as a higher order pays.

\paragraph{Types.} The integer weights must be exact in \code{T}. Every type with 24 significant bits or more
satisfies this, and a narrower type is rejected at the call by a constraint. A derivative above the first, in a type
with an epsilon above $10^{-7}$, keeps few correct digits by every method, and it warns at compile time
(\cref{sec:api-warning}).

\paragraph{The scale.} The default steps assume that $f$ varies on the scale $\max(\abs{x}, 1)$ in the units of $x$.
That is a rule of thumb, not a property of $f$. \Cref{sec:using-pitfalls} shows the symptoms when it fails, and the
cure.

\paragraph{Domain edges.} Ridders' method needs the nearest point at which $f$ is not finite to be a minimum distance
away. The reason is that the levels it spends restarting come out of its budget of \num{\valRiddersLevelCap}. Take
$\log x$ at $x = 2^{-k}$, with the default first-derivative step. The method returned a value down to
$k = \num{\valRiddersEdgeLimitExponent}$, with an estimate grown to flag the loss. But it kept ten correct digits only
down to $k = \num{\valRiddersEdgeTenDigitsExponent}$, and it used its whole budget from
$k = \num{\valRiddersEdgeCapExponent}$ (measured, \cref{fig:ridders-edge}). For ten digits, keep the nearest bad point
at least $2^{-\num{\valRiddersEdgeTenDigitsExponent}} \max(\abs{x}, 1)$ away. The design record derives the first
limit, $2^{-34} \max(\abs{x}, 1)$, from the initial step $s/32$ and the cap. The same derivation gives
$2^{-31} \max(\abs{x}, 1)$ for $f''$, whose rule reaches $h$ from the initial step $s/4$. It gives
$2^{-30} \max(\abs{x}, 1)$ for $f'''$ and $f''''$, whose rules reach $2h$; the record's $2^{-31}$ for these two is a
halving too close. On $\log x$, tested at every half halving of $x$, a value was returned down to
$2^{-\num{\valRiddersEdgeLimitExponentSecond}}$ for $f''$, and down to
$2^{-\num{\valRiddersEdgeLimitExponentThird}}$ and $2^{-\num{\valRiddersEdgeLimitExponentFourth}}$ for $f'''$ and
$f''''$ (measured).

A pole, where $f$ is finite on both sides, behaves otherwise. The restart after an untrusted result now descends as it
does after a non-finite value, so a pole close to $x$ is resolved. The library's test poles lie $10^{-3}$ ahead of $x$
and $10^{-4}$ behind it, far inside the limits above. There, \num{\valPoleCovered} of \num{\valPoleCases} results
for $f'$ to $f''''$ were covered, with errors from \num{\valPoleRelErrorMin} to \num{\valPoleRelErrorMax}
(measured). Closer, a pole does not make the method fail: it returns a wrong value. We swept
\num{\valPoleSweepCases} poles, at $1.25$ to $1.75$ times $2^{-k}$ from $x$ for $k = \num{\valPoleSweepFirstK}$ to
\num{\valPoleSweepLastK}, ahead and behind, for $f'$ to $f''''$. None of the results was \code{non_finite}. Every
one of the \num{\valPoleSweepUncovered} uncovered results was flagged: its estimate was at least a quarter of its
value (measured). Near a pole, test \code{abs_error}.

Exactly at an edge, the one-sided methods apply: \code{Forward<P>} and \code{Backward<P>}, or, with an error
estimate, \code{ForwardRidders} and \code{BackwardRidders} (\cref{sec:ridders-onesided}). They have a reach of their
own. Take $\log x$ at $x = 2^{-k}$. The estimate of \code{ForwardRidders} covered its error up to
$k = \num{\valForwardRiddersReachCoveredExponentFirst}$, \num{\valForwardRiddersReachCoveredExponentSecond},
\num{\valForwardRiddersReachCoveredExponentThird} and \num{\valForwardRiddersReachCoveredExponentFourth} for $f'$
to $f''''$. It flagged the result, with an estimate of at least a quarter of its value, from
$k = \num{\valForwardRiddersReachFlaggedExponentFirst}$, \num{\valForwardRiddersReachFlaggedExponentSecond},
\num{\valForwardRiddersReachFlaggedExponentThird} and \num{\valForwardRiddersReachFlaggedExponentFourth} to the end
of the study, $k = \num{\valForwardRiddersReachLastK}$. Between the covered and the flagged range, $f'$ at
$k = \num{\valForwardRiddersReachSilentExponentFirst}$ was \num{\valForwardRiddersReachSilentRelErrorFirst} off, with
an estimate of \num{\valForwardRiddersReachSilentRelEstimateFirst}, relative. The band between the two is narrow, but
not always a single point. On $\log x$, $\sqrt{x}$ and $\sqrt[3]{x}$ at four points per halving,
\num{\valForwardRiddersBandSilent} of \num{\valForwardRiddersBandCases} results for $f'$ were neither covered nor
flagged (\valForwardRiddersBandFunctions). Each was within \num{\valForwardRiddersBandWorstRatio} times its estimate
(measured).

\paragraph{Aliasing.} A function that oscillates on a scale far below the step can lie at every sample. Samples that
happen to lie on a polynomial fool every sampling method, Ridders' included (\cref{sb:aliasing}).

\paragraph{Noise and kinks.} A function less accurate than $\epsmach$ needs a larger step (\cref{sec:rounding-noise}).
A kink or a jump within a step of $x$ gives a wrong value from every method that samples across it
(\cref{sec:using-pitfalls}).

\paragraph{Overflow.} A final quotient that overflows, as in a fourth derivative with a tiny step, gives
\code{non_finite} at run time. In a constant expression under GCC, it gives a compile error
(\cref{sb:constexpr-overflow}).

\section{Guarantees}\label{sec:guarantees-guarantees}

Each guarantee below is stated with the hypothesis that it rests on, and with how it is verified.
\begin{itemize}
  \item \emph{No exceptions.} The library throws nothing and catches nothing; failures are \code{std::expected}
    values. An exception thrown by \code{fn}, or by the operations of \code{T}, passes through unchanged. Verified by
    inspection: the headers contain no \code{throw} and no \code{try}.
  \item \emph{Where \code{fn} is called.} Only at finite points, since the overflow checks come before any point is
    formed. A stencil calls \code{fn} once per node. \code{Forward<P>} never calls it below $x$, and
    \code{Backward<P>} never above. A step that is not positive and finite returns \code{non_finite} before any call.
    Ridders' method calls \code{fn} at most \num{\valRiddersMaxEvalsFirst}, \num{\valRiddersMaxEvalsSecond},
    \num{\valRiddersMaxEvalsThird} and \num{\valRiddersMaxEvalsFourth} times for $f'$ to $f''''$.
    \code{ForwardRidders} never calls it below $x$, and \code{BackwardRidders} never above, restarts included. They
    call it at most \num{\valForwardRiddersMaxEvalsFirst}, \num{\valForwardRiddersMaxEvalsSecond},
    \num{\valForwardRiddersMaxEvalsThird} and \num{\valForwardRiddersMaxEvalsFourth} times. Verified by tests that
    count and log the calls.
  \item \emph{No shared state.} The headers hold no global, static or thread-local variable other than compile-time
    constants. And \code{fn} is called on the calling thread, one call at a time. So concurrent calls of
    \code{derivative} are as safe as the concurrent calls of \code{fn} that they make. Verified by inspection.
  \item \emph{No allocation.} The library uses only fixed-size arrays and \code{std::optional}, on the stack. A
    \code{T} that allocates memory does so in its own operations. Verified by inspection.
  \item \emph{Failures are values, also at compile time.} Every failure is detected before the arithmetic that would
    produce it. So it is an error value in a constant expression too, with one exception: an overflow of the final
    quotient, which GCC rejects at compile time. Verified by \code{static_assert}s on overflowing, non-finite and
    saturating cases.
  \item \emph{Exact weights.} The integer weights are exact in every type with 24 significant bits or more. The
    largest is 6,222,216, below $2^{24}$. A larger value would make the generator's constant evaluation fail, and a
    narrower type is rejected by a constraint of \code{derivative}. Verified at every compilation.
  \item \emph{Exact on polynomials.} Take a dyadic $x$, that is, an integer times a power of two, and a power-of-two
    $h$. Then a stencil is exact for polynomials of degree up to $\dorder + \aorder - 1$, provided that every
    intermediate result is exact in the type. Verified by \code{static_assert}s that also check the bit budget of
    each case.
  \item \emph{Compile time equals run time.} The same call gives the same bits at compile time and at run time,
    provided that the compiler neither contracts \code{a * b + c} into a fused multiply-add nor uses extended
    precision. The build pins neither. Verified on the five toolchains of the continuous integration (GCC, Clang,
    MSVC, clang-cl, Emscripten): by tests that compare constant and run-time results, and by
    \code{code/07-using-constexpr.cpp}.
  \item \emph{\code{root_epsilon}} is within an ulp of $\epsmach^{1/N}$. Verified against references computed to 80
    digits, for selected N from 2 to 2000, in \code{float} and \code{double}; and against 60-digit references for the
    multiprecision types.
  \item \emph{Types.} \code{derivative} is tested on \code{float}, \code{double}, \code{cpp_bin_float_50} with and
    without expression templates, and \code{cpp_dec_float_50}. The \code{float} derivatives above the first are
    tested with the warning silenced at the call. \code{long double} satisfies the concept, and is tested in
    \code{root_epsilon} only.
\end{itemize}
Three things are not guaranteed: the accuracy, which is a model and a measurement (\cref{sec:guarantees-accuracy});
Ridders' error estimate, which is an estimate (\cref{sec:guarantees-ridders}); and the default step, which is a rule
of thumb.

\section{How it is verified}\label{sec:guarantees-verification}

The tests of \code{nxx::deriv} are part of one doctest executable, a single test program, and there are four build
tests, which check what a compilation reports. \Cref{tab:tests} lists them. Their reference values for the derivatives come from closed forms,
from the published tables of Fornberg\cite{fornberg1988} and Chapra (Chapter 23),\cite{chapra2015} or from an
independent prototype, never from \nxx{} itself. The ceilings on \code{abs_error} for the multiprecision types are
different: they are twice the estimates of the first run. They are regression baselines, values recorded from a
run to catch later changes, and \code{docs/deriv/design.md} records them as such. They are not independent
references. Each tolerance is justified by an error estimate in a comment. The stencils, for example, are held to
$(\errc + 2\ampl)\,\epsmach^{P/(P+N)}$ relative on $\mathrm{e}^x$, for a selection of rules and orders, two-node
rules included. That tolerance is the truncation error plus the rounding of the samples and of the points. The tests'
comment calls it an estimate, not a bound, and reports measured errors 4 to 16 times below it on its three points.
The tests run on the five toolchains of the continuous integration, with warnings as errors.

% Between the paragraph on the tests and the one on the reviews (the language pass, PLAN.md 12.9): at the head
% of the section it floated to the foot of the page and split the paragraph; after the whole paragraph it
% left a third of a page empty.
\begin{table}[htb]
\sidecaption{The tests of \code{nxx::deriv}, in \code{tests/}.\label{tab:tests}}
\centering\footnotesize
\begin{tabular}{@{}>{\raggedright}p{34mm}>{\raggedright\arraybackslash}p{75mm}@{}}
\toprule
File & What it checks\\
\midrule
\code{test_math.cpp} & the concept \code{real} on accepted and rejected types; \code{abs}, \code{isfinite};
  \code{root_epsilon} against 80-digit roots, the same at compile time and run time\\
\code{test_deriv.cpp} & \code{Central<2>} with the default step on float, double and multiprecision types; two calls
  of \code{fn}; non-finite results as errors\\
\code{test_stencils.cpp} & stencils of each family with their default step against $(\errc + 2\ampl)\,
  \epsmach^{P/(P+1)}$; one-sided rules never cross $x$; one call per node; failures; multiprecision\\

\code{test_callables.cpp} & \code{fn} called through \code{std::invoke} with a \code{T const} lvalue: member pointers,
  callables that reject rvalues\\
\code{test_ridders.cpp} & exact results on polynomials with the floor as the estimate; accuracy and coverage
  against closed forms and the prototype; restarts; failures, untrusted results and the level cap; multiprecision\\
\code{test_orders.cpp} & orders 1 to 4: the generated tables against the moment conditions and published rows;
  exactness on polynomials; steps outside the range of $h^N$; overflow; Ridders' evaluations and accuracy, poles
  close to $x$\\
\code{test_onesided.cpp} & \code{ForwardRidders} and \code{BackwardRidders}: exact on polynomials; accuracy and
  coverage against closed forms and the prototype; domain edges, a pole behind $x$, a kink at $x$; \code{fn} called on
  one side only, restarts included; the mirror; failures and cost; multiprecision\\
\code{precision_warning.cpp} & four builds: a float $f''$ warns through each overload, and a pragma or the macro
  silences it\\
\bottomrule
\end{tabular}
\end{table}

Each phase was reviewed by specialised review agents, and Phases~4 and~5 also by another vendor's model (OpenAI
Codex). The findings and decisions are recorded in the design record (\code{docs/deriv/PHASE3_API.md} to \code{docs/deriv/PHASE5_API.md},
\code{docs/deriv/design.md}). The book's own examples build and run as tests. And each time it runs, the book's data
generator checks every table against its moment conditions, and the nitrogen function against an independent exact
evaluation.
